
\documentclass{article}
\usepackage{colortbl}
\usepackage{makecell}
\usepackage{multirow}
\usepackage{supertabular}

\begin{document}

\newcounter{utterance}

\centering \large Interaction Transcript for game `eqbench', experiment `eq\_bench', episode 121 with reflexion{-}llama{-}3.1{-}8b{-}instruct.
\vspace{24pt}

{ \footnotesize  \setcounter{utterance}{1}
\setlength{\tabcolsep}{0pt}
\begin{supertabular}{c@{$\;$}|p{.15\linewidth}@{}p{.15\linewidth}p{.15\linewidth}p{.15\linewidth}p{.15\linewidth}p{.15\linewidth}}
    \# & \multicolumn{2}{c}{Player} && \multicolumn{2}{c}{Game Master} \\
    \hline

    \theutterance \stepcounter{utterance}  
    & & & \multicolumn{4}{p{0.6\linewidth}}{
        \cellcolor[rgb]{0.9,0.9,0.9}{
            \makecell[{{p{\linewidth}}}]{
                \texttt{\tiny{[P1$\langle$GM]}}
                \texttt{You have attempted this same task before and did not succeed. These are the plans you wrote after those attempts.} \\
\texttt{Trial \#1: New plan: I got stuck in a loop of analyzing the dialogue without considering the underlying dynamics of the conversation. I focused on individual statements without taking into account the emotional tone and the power dynamics between Andrew and Sonya. Next time, I will pay closer attention to the emotional cues and the language used by each character, particularly the use of rhetorical questions, metaphors, and emotional appeals. I will also consider the context of the conversation, including the stakes and the relationships between the characters. By doing so, I will be able to better understand the emotional responses of the characters and make more accurate predictions.} \\
\\ 
\texttt{Trial \#2: New plan: I got stuck in a loop of analyzing the dialogue without considering the underlying dynamics of the conversation. I focused on individual statements without taking into account the emotional tone and the power dynamics between Andrew and Sonya. Next time, I will pay closer attention to the emotional cues and the language used by each character, particularly the use of rhetorical questions, metaphors, and emotional appeals. I will also consider the context of the conversation, including the stakes and the relationships between the characters. By doing so, I will be able to better understand the emotional responses of the characters and make more accurate predictions.} \\
\\ 
\texttt{To achieve this, I will:} \\
\\ 
\texttt{1. Identify the emotional tone of each character's statements, including the use of emotional appeals, rhetorical questions, and metaphors.} \\
\texttt{2. Analyze the power dynamics between Andrew and Sonya, including their roles, relationships, and the stakes of the conversation.} \\
\texttt{3. Consider the context of the conversation, including the setting, the topic, and the goals of each character.} \\
\texttt{4. Use this information to make more accurate predictions about the emotional responses of the characters, particularly Andrew.} \\
\\ 
\texttt{By following this plan, I will be able to better understand the underlying dynamics of the conversation and make more accurate predictions about the emotional responses of the characters.} \\
\\ 
\\ 
\\ 
\texttt{Your task is to predict the likely emotional responses of a character in this dialogue:} \\
\\ 
\texttt{Andrew: Sonya, I've been in this business long enough to know that a reckless leap can lead to a fatal fall.} \\
\texttt{Sonya: But Andrew, playing it safe won't get us anywhere. We need to take risks to innovate.} \\
\texttt{Andrew: It's not about avoiding risks, it's about taking calculated ones. You're playing with fire here.} \\
\texttt{Sonya: And what if I'm ready to get burned? What if that's what it takes to light up the path?} \\
\texttt{Andrew: You're not just risking your own future, Sonya. There are employees, investors, stakeholders involved.} \\
\texttt{Sonya: I'm well aware of that, Andrew. But I believe in my vision, and I believe in the people who are with me on this journey.} \\
\texttt{Andrew: It's not about belief, Sonya. It's about strategy, it's about calculated moves.} \\
\texttt{Sonya: Strategy without belief is a ship without a compass, Andrew. It might keep you afloat, but it won't get you anywhere.} \\
\texttt{Andrew: And belief without strategy is a ship without a rudder. You'll be lost at sea.} \\
\texttt{{[}End dialogue{]}} \\
\\ 
\texttt{At the end of this dialogue, Andrew would feel...} \\
\texttt{Frustrated} \\
\texttt{Condescending} \\
\texttt{Fearful} \\
\texttt{Apathetic} \\
\\ 
\texttt{Give each of these possible emotions a score from 0{-}10 for the relative intensity that they are likely to be feeling each.} \\
\\ 
\texttt{You must output in the following format, including headings (of course, you should give your own scores), with no additional commentary:} \\
\\ 
\texttt{Frustrated: <score>} \\
\texttt{Condescending: <score>} \\
\texttt{Fearful: <score>} \\
\texttt{Apathetic: <score>} \\
\\ 
\\ 
\texttt{{[}End of answer{]}} \\
\\ 
\texttt{Remember: zero is a valid score, meaning they are likely not feeling that emotion. You must score at least one emotion > 0.} \\
\\ 
\texttt{Your answer:} \\
            }
        }
    }
     \\ \\

    \theutterance \stepcounter{utterance}  
    & \multicolumn{4}{p{0.6\linewidth}}{
        \cellcolor[rgb]{0.561,0.737,0.561}{
            \makecell[{{p{\linewidth}}}]{
                \texttt{\tiny{[P1$\rangle$GM]}}
                \texttt{Frustrated: 8} \\
\texttt{Condescending: 2} \\
\texttt{Fearful: 0} \\
\texttt{Apathetic: 0} \\
            }
        }
    }
    & & \\ \\

    \theutterance \stepcounter{utterance}  
    & & & \multicolumn{2}{p{0.3\linewidth}}{
        \cellcolor[rgb]{0.9,0.9,0.9}{
            \makecell[{{p{\linewidth}}}]{
                \texttt{\tiny{[GM$|$GM]}}
                \texttt{Frustrated: 8} \\
\texttt{Condescending: 2} \\
\texttt{Fearful: 0} \\
\texttt{Apathetic: 0} \\
            }
        }
    }
    & & \\ \\

    \theutterance \stepcounter{utterance}  
    & & & \multicolumn{2}{p{0.3\linewidth}}{
        \cellcolor[rgb]{0.9,0.9,0.9}{
            \makecell[{{p{\linewidth}}}]{
                \texttt{\tiny{[GM$|$GM]}}
                \texttt{Frustrated: 8} \\
\texttt{Condescending: 4} \\
\texttt{Fearful: 6} \\
\texttt{Apathetic: 0} \\
            }
        }
    }
    & & \\ \\

    \theutterance \stepcounter{utterance}  
    & & & \multicolumn{2}{p{0.3\linewidth}}{
        \cellcolor[rgb]{0.9,0.9,0.9}{
            \makecell[{{p{\linewidth}}}]{
                \texttt{\tiny{[GM$|$GM]}}
                \texttt{game\_result = LOSE} \\
            }
        }
    }
    & & \\ \\

    \theutterance \stepcounter{utterance}  
    & & \multicolumn{4}{p{0.6\linewidth}}{
        \cellcolor[rgb]{0.95,0.95,0.95}{
            \makecell[{{p{\linewidth}}}]{
                \texttt{\tiny{[P1$\leftrightarrow$P]}}
                \texttt{Post{-}episode Reflexion (not sent to GM):} \\
\texttt{New plan: I got stuck in a loop of analyzing the dialogue without considering the underlying dynamics of the conversation. I focused on individual statements without taking into account the emotional tone and the power dynamics between Andrew and Sonya. Next time, I will pay closer attention to the emotional cues and the language used by each character, particularly the use of rhetorical questions, metaphors, and emotional appeals. I will also consider the context of the conversation, including the stakes and the relationships between the characters. By doing so, I will be able to better understand the emotional responses of the characters and make more accurate predictions.} \\
\\ 
\texttt{However, upon reviewing my previous attempts, I realize that I provided the same plan twice. To improve my strategy, I will add a new layer of analysis by considering the following:} \\
\\ 
\texttt{1. Identify the emotional tone of each character's statements, including the use of emotional appeals, rhetorical questions, and metaphors.} \\
\texttt{2. Analyze the power dynamics between Andrew and Sonya, including their roles, relationships, and the stakes of the conversation.} \\
\texttt{3. Consider the context of the conversation, including the setting, the topic, and the goals of each character.} \\
\texttt{4. Use this information to make more accurate predictions about the emotional responses of the characters, particularly Andrew.} \\
\\ 
\texttt{Additionally, I will also consider the following:} \\
\\ 
\texttt{5. Look for moments where Andrew and Sonya use language that reveals their underlying emotions, such as phrases that indicate frustration, concern, or excitement.} \\
\texttt{6. Analyze the way Andrew and Sonya respond to each other's statements, including their tone, body language, and verbal cues.} \\
\texttt{7. Consider the implications of Sonya's statement "And what if I'm ready to get burned? What if that's what it takes to light up the path?" and how it might affect Andrew's emotional response.} \\
\\ 
\texttt{By incorporating these additional layers of analysis, I will be able to develop a more nuanced understanding of the emotional dynamics at play in the conversation and make more accurate predictions about Andrew's emotional response.} \\
            }
        }
    }
    & \\ \\

\end{supertabular}
}

\end{document}
